%% toc: P15 164 164.6
\chapter{Carnap 1950: the infinite system, interpretations, and classical laws}
\label{sec:p15-w02b-p15-029}

This chapter continues
Chapters~\ref{sec:p15-carnap-logical-probability-and-induction}
and~\ref{sec:p15-carnap-confirmation-relevance-and-reliability}. Its
purpose is to supply the parts of the 1950 first edition of
\emph{Logical Foundations of Probability} that those chapters leave as
boundaries: the infinite system \(\mathfrak L_\infty\) and its limit
functions, the three explanations of probability\(_1\) in sec.~41, the
binomial law and Bernoulli's theorem of secs.~95--96, and the
variety-of-instances requirement of sec.~110I. A source-to-host bridge
records, item by item, whether Carnap printed a proof or the proof given
here is an independent reconstruction. A final part reconstructs the
first five sections of Zabell's historical chapter. The inputs are the
finite systems \(\mathfrak L_N\), regular and symmetric measures,
and \(m^*\), \(c^*\) as defined in
\cref{def:p15-carnap-state-range,def:p15-carnap-regular-confirmation,def:p15-carnap-star}.
The 1952 \emph{Continuum of Inductive Methods} and the 1962 second
edition are not used.

\section{The infinite system and its limit functions}
\label{sec:p15-w02b-p15-029-infinite}

\begin{hbdefinition}{def:p15-w02b-p15-029-systems}{The systems \(\mathfrak L_N\) and \(\mathfrak L_\infty\)}
The system \(\mathfrak L_\infty\) has a denumerable sequence of
individual constants \(a_1,a_2,\ldots\), designating all individuals of
its domain, with distinct constants designating distinct individuals.
It has finitely many primitive predicates, the identity sign \(=\), the
connectives, the universal quantifier and individual variables.
For \(N\geq1\), \(\mathfrak L_N\) has only \(a_1,\ldots,a_N\) and
otherwise the same signs. A sentence \(i\) of \(\mathfrak L_\infty\)
\emph{occurs} in \(\mathfrak L_N\) when every constant in \(i\)
has index at most \(N\). In each system a state description decides
every atomic sentence; an identity sentence with two different constants
is L-false, and one with the same constant twice is L-true. A sentence
is \emph{general} if it contains a variable, \emph{nongeneral}
otherwise, and \emph{purely general} if it is general and contains no
individual constant \cite[sec.~15A; D16-6f--h]{Carnap1950Probability}.
\end{hbdefinition}

\begin{hbdefinition}{def:p15-w02b-p15-029-limits}{Limit functions: D56-1 and D56-2}
Let \(m_1,m_2,\ldots\) be regular \(m\)-functions for
\(\mathfrak L_1,\mathfrak L_2,\ldots\), and let \(c_N\) be the
regular \(c\)-function based on \(m_N\). For a sentence \(i\) of
\(\mathfrak L_\infty\), put
\[
 m_\infty(i)\defeq\lim_{N\to\infty}m_N(i)
\]
when this limit exists, the sequence being taken over the systems in
which \(i\) occurs; otherwise \(m_\infty\) has no value for \(i\).
Likewise \(c_\infty(h,e)\defeq\lim_N c_N(h,e)\) when the limit
exists, taken over those \(N\) for which \(c_N(h,e)\) is defined.
These are D56-1 and D56-2. By T56-2, \(c_\infty(h,e)\) has a value
exactly when \(e\) is not L-false in infinitely many
\(\mathfrak L_N\) and the values \(c_N(h,e)\) for those \(N\)
converge \cite[sec.~56, pp.~302--304]{Carnap1950Probability}.
\end{hbdefinition}

Two hypotheses are built into these definitions. The domain is
countable and every individual is named, so that \(\mathfrak L_\infty\)
is approached by the finite systems obtained by deleting all but the
first \(N\) names. Second, a value is defined only when the limit
exists; the definition does not itself assert existence. Carnap rejects
the alternative of a measure on the state descriptions of
\(\mathfrak L_\infty\) itself: that set has the cardinality of the
continuum, and a direct quotient \(m_\infty(e\land h)/m_\infty(e)\)
gives no value when \(m_\infty(e)=0\), although the finite quotients
may still converge. For \(e=\forall x\,M(x)\) with \(M\) factual,
\(c_N(M(a_1),e)=1\) for every \(N\), so \(c^*_\infty(M(a_1),e)=1\),
although \(m^*_\infty(e)=0\) by \cref{thm:p15-w02b-p15-029-monadic-limit}
\cite[sec.~56, discussion following T56-2; T56-4]{Carnap1950Probability}.

In the appendix the existence of the \(m^*\)-limit is asserted for
every general sentence, together with an effective and practicable
reduction to a nongeneral sentence, but no proof is printed. That
passage stands under the declared restriction of the remainder of
sec.~110 to systems whose primitive predicates all have degree one
\cite[sec.~110A, pp.~566--567]{Carnap1950Probability}. The next two
results supply proofs of those claims for the monadic systems.
Throughout, \(\mathfrak L_N\) and \(\mathfrak L_\infty\) have
\(\pi\geq1\) primitive predicates, all of degree one, and
\(\kappa=2^\pi\) \(Q\)-properties \(Q_1,\ldots,Q_\kappa\). The
quantifier rank of a formula is its maximal nesting depth of quantifiers.

\begin{hblemma}{thm:p15-w02b-p15-029-saturation}{proved=annex:p15-w02b-p15-029-saturation}{Counts above the quantifier rank do not matter}
Let \(i\) be a sentence of a monadic system whose individual constants
are exactly \(b_1,\ldots,b_m\) (\(m\geq0\)), and let \(r\) be its
quantifier rank. Let \(\mathcal A,\mathcal B\) be structures for the
predicates and these constants, with \(b_1,\ldots,b_m\) denoting
distinct elements in each. Suppose that each \(b_l\) has the same
\(Q\)-property in \(\mathcal A\) and \(\mathcal B\), and that for
each \(j\) the numbers of elements of \(Q_j\) not denoted by any
\(b_l\) are either equal in \(\mathcal A\) and \(\mathcal B\) or are
both at least \(r\) (possibly infinite). Then
\(\mathcal A\models i\) if and only if \(\mathcal B\models i\).
\end{hblemma}

For a word \(u=(u_1,\ldots,u_m)\) of \(Q\)-properties, let
\(D_{r,u}\) be the finite structure with elements \(b_1,\ldots,b_m\)
of \(Q\)-properties \(u_1,\ldots,u_m\) and exactly \(r\) further
elements in every \(Q_j\). Let \(\langle u\rangle\) be the nongeneral
conjunction \(Q_{u_1}(b_1)\land\cdots\land Q_{u_m}(b_m)\).

\begin{hbtheorem}{thm:p15-w02b-p15-029-monadic-limit}{proved=annex:p15-w02b-p15-029-monadic-limit}{Existence, effective reduction, and 110A (5)}
Let \(i\) be any sentence of monadic \(\mathfrak L_\infty\), with
constants \(b_1,\ldots,b_m\) and quantifier rank \(r\). Put
\[
 S_i\defeq\{u\in\{1,\ldots,\kappa\}^m:D_{r,u}\models i\},\qquad
 i'\defeq\bigvee_{u\in S_i}\langle u\rangle ,
\]
with the empty disjunction read as \(\neg t\) and, when \(m=0\), the
single empty word giving \(t\) or \(\neg t\).
\begin{enumerate}
 \item \(m^*_\infty(i)\) exists and equals \(m^*(i')\). The sentence
       \(i'\) is nongeneral, contains no individual constant not in
       \(i\), and is computed by evaluating \(i\) in the
       \(\kappa^m\) finite structures \(D_{r,u}\).
 \item If \(i\) is purely general, \(m^*_\infty(i)\) is \(0\) or
       \(1\); it is \(1\) exactly when \(D_{r,\varnothing}\models i\).
 \item If \(M\) is a molecular predicate of width \(0<w<\kappa\),
       \(m^*_\infty(\forall x\,M(x))=0\).
 \item If \(m^*_\infty(e)>0\), then \(c^*_\infty(h,e)\) exists for
       every \(h\) and equals
       \(m^*_\infty(e\land h)/m^*_\infty(e)\), computable by (1).
\end{enumerate}
\end{hbtheorem}

Item (2) is Carnap's 110A (5); item (3) is the remark following it,
which leads to his formula (12), \(c^*(l,e)=0\) in
\(\mathfrak L_\infty\). The procedure in (1) is effective. Its cost
is \(\kappa^m\) evaluations of \(i\) in structures with
\(m+r\kappa\) elements; no claim is made here that this agrees with
the practicable procedure Carnap announces for his second volume. The
reduced sentence \(i'\) need not be L-equivalent to \(i\): for
\(i=\exists x\,Q_1(x)\), \(i'=t\), although \(i\) is not L-true.

\begin{hbexample}{def:p15-w02b-p15-029-limit-example}{A general sentence with a constant}
Take \(\pi=1\), so \(\kappa=2\), and
\(i=P(a_1)\land\exists x(x\ne a_1\land\neg P(x))\), of rank one.
\(D_{1,u}\) contains \(a_1\) and one further element in each of
\(P,\neg P\); so \(S_i=\{P\}\), \(i'=P(a_1)\), and
\(m^*_\infty(i)=1/2\). Directly, \(i\) fails in a state where
\(a_1\) is \(P\) exactly when every individual is \(P\), and that
state has \(m^*_N\)-mass \(1/(N+1)\); so
\(m^*_N(i)=1/2-1/(N+1)\), which tends to \(1/2\).
\end{hbexample}

The comparative theorem T81-2a in the 1950 text transfers the finite
adequacy result of \cref{thm:p15-carnap-comparative} to
\(\mathfrak L_\infty\). Its printed proof rests on one logical step:
L-truth of the conditions of D81-1c in \(\mathfrak L_\infty\) passes to
every \(\mathfrak L_N\) of a final segment (T20-11 is cited for it)
\cite[sec.~81, pp.~440--441]{Carnap1950Probability}. The following
results prove that step for arbitrary general sentences and then the
transfer.

\begin{hbtheorem}{thm:p15-w02b-p15-029-ltruth-transfer}{proved=annex:p15-w02b-p15-029-ltruth-transfer}{L-truth passes from the infinite system to a final segment}
Let \(i\) be a sentence of \(\mathfrak L_\infty\).
\begin{enumerate}
 \item (Monadic systems, no choice principle.) Let \(i\) have
       \(m\) constants, largest index \(n_0\), and rank \(r\geq1\).
       If \(i\) is L-true in \(\mathfrak L_\infty\), then \(i\) is
       L-true in every \(\mathfrak L_N\) with
       \(N\geq\max\{n_0,\,m+\kappa(r-1)+1\}\). Conversely, if
       \(i\) is L-true in every \(\mathfrak L_N\) with
       \(N\geq\max\{n_0,\,m+\kappa r\}\), then it is L-true in
       \(\mathfrak L_\infty\).
 \item (Arbitrary finite relational signature, assuming the Axiom of
       Choice as in \cref{thm:fol-compactness,thm:downward-ls}.)
       If \(i\) is L-true in \(\mathfrak L_\infty\), there is
       \(N_0\) such that \(i\) is L-true in every
       \(\mathfrak L_N\) with \(N\geq N_0\).
\end{enumerate}
\end{hbtheorem}

\begin{hbtheorem}{thm:p15-w02b-p15-029-t81-transfer}{proved=annex:p15-w02b-p15-029-t81-transfer}{T81-2a for general sentences}
Let \(h,e,h',e'\) be sentences of \(\mathfrak L_\infty\), general or
not, with \(e,e'\) not L-false in any system in which they occur, and
suppose \(\operatorname{MC}(h,e;h',e')\) holds in
\(\mathfrak L_\infty\) (\cref{def:p15-carnap-comparative}). Then
\(\operatorname{MC}(h,e;h',e')\) holds in every \(\mathfrak L_N\)
of a final segment, so \(c_N(h,e)\geq c_N(h',e')\) there for every
regular sequence. Hence
\(c_\infty(h,e)\geq c_\infty(h',e')\) whenever both limits exist.
In the monadic systems the final segment is explicit by
\cref{thm:p15-w02b-p15-029-ltruth-transfer}, and for \(c^*\) both
limits exist whenever \(m^*_\infty(e)>0\) and \(m^*_\infty(e')>0\).
\end{hbtheorem}

The limit-existence hypothesis cannot be dropped from the conclusion in
general; it is the condition T81-2a itself states. The converse, T81-2b,
is stated by Carnap for nongeneral sentences only, and is not extended
here.

\section{Probability\texorpdfstring{\(_1\)}{1} as support, betting quotient, and estimate}
\label{sec:p15-w02b-p15-029-interpretations}

Section~41 clarifies the explicandum before any explicatum is
defined. It gives three explanations of the probability\(_1\) of
\(h\) on \(e\): a measure of evidential support, a fair betting
quotient, and an estimate of relative frequency
\cite[sec.~41, pp.~162--175]{Carnap1950Probability}. The circularity
among them is permitted, Carnap says, only because these are
clarifications of explicanda; the later system defines its concepts
in a noncircular chain.

\begin{hbdefinition}{def:p15-w02b-p15-029-betting}{Stakes, betting quotient, estimate}
In a bet on \(h\) between \(X_1\) and \(X_2\), \(X_1\) receives
\(u_2\) if \(h\) is true and pays \(u_1\) if it is false;
\(u_1,u_2\geq0\) and \(u_1+u_2>0\). The \emph{betting ratio} is
\(u_1:u_2\) and the \emph{betting quotient} is
\(q=u_1/(u_1+u_2)\). The bettors pool their knowledge \(e\). The
statement that the probability\(_1\) of \(h\) on \(e\) is \(q\)
is read as: a bet on \(h\) with quotient \(q\) is \emph{fair}, that
is, favours neither side. If a magnitude has finitely many possible
values \(u_1,\ldots,u_n\), its \emph{estimate} on \(e\) is the
probability\(_1\)-weighted mean of these values (41D (3)).
\end{hbdefinition}

\begin{hbtheorem}{thm:p15-w02b-p15-029-betting}{proved=annex:p15-w02b-p15-029-betting}{The arithmetic of sec.~41}
\begin{enumerate}
 \item Every betting quotient lies in \([0,1]\), and \(q\) for
       \(h\) is the quotient \(1-q\) for \(\neg h\) from the other
       side (41B (1), (2)).
 \item If \(X_1\) makes the same bet with quotient \(q\) on each of
       \(n\) individuals of a class \(K\) having \(M\), and the
       relative frequency of \(M\) in \(K\) is \(r\), his total gain
       is \(n(u_1+u_2)(r-q)\).
 \item If the possible numbers of \(M\)-members of \(K\) carry
       probabilities\(_1\) summing to one, the estimates satisfy
       \(r'=m'/n\) and \(g'=n(u_1+u_2)(r'-q)\) (41D (4), (5)). The
       estimated gain is zero exactly when \(q=r'\) (41D (6)).
 \item (41A) Let support values be nonnegative, at most one, and
       additive over L-exclusive hypotheses, with the L-true sentence
       at one. If \(j_1,\ldots,j_n\) are L-exclusive in pairs and
       L-disjunct and equally supported by \(e\), each has support
       \(1/n\) and a disjunction of \(k\) of them has \(k/n\).
\end{enumerate}
\end{hbtheorem}

Carnap regards the assumptions in item (4) as arbitrary when only
evidential support is in view; they are justified by the betting and
estimate explanations, which yield items (1)--(3). When \(e\) states
that \(r\) is the frequency of \(M\) in \(K\), the only quotient
neither bettor will refuse is \(r\) (41C); when \(r\) is unknown, the
fair quotient is its estimate (41D (6)). Two further claims make the
estimate explanation coherent: the estimate of the frequency of \(M\)
in an unobserved class does not depend on the class (41D (7)), and it
equals the probability\(_1\) of the singular prediction (41D (8)).
Carnap prints the derivation of (8) from (7) and defers (7) to the
later T106-1d. More generally, among hypotheses of equal
probability\(_1\) \(q\), the estimated relative frequency of truth
is \(q\) (41D (9), later T104-2c). For symmetric \(c\)-functions these
are the identities of \cref{thm:p15-carnap-frequency-reliability}.

The probability\(_2\) of an outcome is a physical magnitude whose value
may be unknown, and one may ask for its probability\(_1\) on evidence.
A probability\(_1\) statement is analytic, so its own probability\(_1\)
on any evidence is \(0\) or \(1\). Carnap applies this to the die
example of 41D: from a statement (ii) that the probability\(_1\) of an
ace is \(1/6\) on symmetric-shape evidence, the factual statement that
the long-run frequency \emph{will be} \(1/6\) does not follow; the
statement that its \emph{estimate} is \(1/6\) does, and a high
probability\(_1\) for a frequency near \(1/6\) follows only with the
restricting conditions of Bernoulli's theorem
(\cref{thm:p15-w02b-p15-029-bernoulli}).

\section{The binomial law and Bernoulli's theorem}
\label{sec:p15-w02b-p15-029-bernoulli}

Here, as in T94-1, a division \(M_1,\ldots,M_p\) is given; \(e\)
states that a population of \(n\) named individuals has \(n_i\)
members of \(M_i\), so \(r_i=n_i/n\); \(h\) is an individual
distribution assigning categories to \(s\) specified members, with
counts \(s_i\leq n_i\); and \(h_{st}\) is the corresponding statistical
distribution. The function \(c\) is any symmetric \(c\)-function, and
the exact value of \(c(h_{st},e)\) is the hypergeometric value of
\cref{thm:p15-carnap-direct}.

\begin{hbtheorem}{thm:p15-w02b-p15-029-binomial}{proved=annex:p15-w02b-p15-029-binomial}{T95-1: the binomial law and its two senses}
\begin{enumerate}
 \item[(A)] For every population,
       \[
        \prod_i\Bigl(1-\frac{s_i-1}{n_i}\Bigr)^{s_i}
        \leq\frac{c(h,e)}{\prod_i r_i^{s_i}}
        \leq\Bigl(1-\frac{s-1}{n}\Bigr)^{-s},
       \]
       (factors with \(s_i=0\) read as one, and \(n\geq s\)). Hence
       \(c(h,e)\) is approximately \(\prod_ir_i^{s_i}\) when \(n\) is
       large relative to \(s\) and each \(n_i\) relative to \(s_i\).
       In the two-category case the law of the sample count differs from
       the binomial law by at most \(s(s-1)/(2n)\) in total variation.
 \item[(B)] If \(n\to\infty\) with \(n_i/n\to r_i\) for every
       \(i\) and \(s,s_i\) fixed, then
       \(c(h,e)\to r_1^{s_1}\cdots r_p^{s_p}\) (T95-1a),
       \(c(h_{st},e)\to\frac{s!}{s_1!\cdots s_p!}r_1^{s_1}\cdots r_p^{s_p}\)
       (T95-1b), and for \(p=2\) this is the binomial law
       \(\binom s{s_1}r_1^{s_1}r_2^{s-s_1}\) (T95-1c).
 \item[(d)] For \(p=2\) and \(0<r_1<1\), put \(q=(s+1)r_1\). The
       binomial values have a unique maximum at \(s_1=\lfloor q\rfloor\)
       when \(q\) is not an integer, and two equal maxima at \(q-1\)
       and \(q\) when it is. In particular the maximum is at \(sr_1\)
       when \(sr_1\) is an integer (T95-1d).
\end{enumerate}
\end{hbtheorem}

\begin{hbtheorem}{thm:p15-w02b-p15-029-bernoulli}{proved=annex:p15-w02b-p15-029-bernoulli}{T96-1: Bernoulli's theorem under its restricting conditions}
Let \(p=2\), \(M=M_1\), \(0<r_1<1\), \(r_2=1-r_1\), and
\(\sigma=\sqrt{sr_1r_2}\).
\begin{enumerate}
 \item[(e)] (Limit theorem, exact form.) For every symmetric \(c\),
       every finite population with \(n\geq s\), and every \(g>0\),
       \[
        c\Bigl(\Bigl|\frac{s_1}s-r_1\Bigr|\leq g,\;e\Bigr)
        \geq1-\frac{r_1r_2}{sg^2}\cdot\frac{n-s}{n-1}
        \geq1-\frac1{4sg^2}.
       \]
       Thus for given \(g,\varepsilon>0\) there is \(s'\) such that the
       value exceeds \(1-\varepsilon\) for every \(s\geq s'\) and every
       population size, which is T96-1e.
 \item[(a)--(d)] (Normal law.) For the binomial values \(B_s\) of
       \cref{thm:p15-w02b-p15-029-binomial}(B), as \(s\to\infty\) with
       \(r_1\) fixed, \(\sigma B_s(s_1)-\phi(\delta/\sigma)\to0\)
       uniformly for \(|\delta|\leq A\sigma\), where
       \(\delta=s_1-sr_1\) and \(0<A<\infty\). For fixed
       \(u_1<u_2\) and \(\delta_l=u_l\sigma\), the \(B_s\)-value of
       \(\{\delta_1\leq\delta\leq\delta_2\}\) tends to
       \(\Phi(u_2)-\Phi(u_1)\), and so does the value of the interval
       widened by \(\tfrac12\) at each end. Hence the symmetric-interval
       value tends to \(2\Phi(u)-1\) and the outside-interval value to
       \(2(1-\Phi(u))\). Along populations with \(n_1/n=r_1\), the
       same interval limits hold for the exact finite-population values
       when \(s^2/n\to0\), and the local limit when
       \(s^{5/2}/n\to0\).
\end{enumerate}
\end{hbtheorem}

Carnap's statements (a)--(d) are approximations: they hold, in his
words, when \(s\) and even \(sr_1r_2\) are sufficiently large and the
population is very large relative to the sample. He gives no proof and
refers to textbooks \cite[sec.~96, pp.~503--505]{Carnap1950Probability}.
The limits above are the exact content that such approximations can
have; Carnap's ranking of the forms (1), (2), (3) by accuracy is a
numerical observation, not a theorem with an error bound, and it is not
asserted here. His statement (e) holds for a fixed \(r_1\); the bound
proved above is uniform in the population and needs no normal
approximation.

\begin{hbexample}{def:p15-w02b-p15-029-chicago}{The source's Chicago example}
Carnap takes \(s=2400\), \(n=3{,}000{,}000\), \(r_1=0.6\), so
\(sr_1r_2=576\), \(\sigma=24\), and \(sr_1=1440\). For the interval
\(1440\pm24\), (c)(1) gives \(u=24.5/24\) and about \(0.692\); the
rougher form gives \(u=1\) and about \(0.682\)
\cite[sec.~96, pp.~506--507]{Carnap1950Probability}. Here
\(s(s-1)/(2n)\) is about \(0.96\), so the total-variation bound of
\cref{thm:p15-w02b-p15-029-binomial}(A) does not by itself certify the
transfer of the binomial value to this finite population, while the
exact bound of (e) gives only about \(0.0008\) for \(g=0.01\), the
relative width of that interval: it is uninformative there. Both comparisons show how much the stated
conditions carry. For \(g=0.05\) the exact bound gives at least
\(0.96\).
\end{hbexample}

\begin{hbwarning}{def:p15-w02b-p15-029-binomial-misuse}{Restricting conditions}
The binomial law in T95-1c presupposes evidence stating the population
frequency \(r_1\). Carnap rejects the traditional use that replaces it
by the frequency observed in an earlier sample, for example
\(\binom{100}{20}(0.12)^{20}(0.88)^{80}\) for twenty aces in the next
hundred throws after twenty-four aces in two hundred: that inference
is predictive, not direct, and the binomial form holds there only
approximately under stronger conditions. Sampling with replacement is,
strictly, sampling from a series of populations with the same counts.
For an infinite population the law holds exactly when \(e\) states the
limit of the frequency, but such evidence can hardly be the observational
evidence actually available \cite[sec.~95, pp.~499--502]{Carnap1950Probability}.
For \(c^*\), the predictive value is given by
\cref{thm:p15-carnap-predictive}, not by the binomial formula.
\end{hbwarning}

\section{The variety of instances}
\label{sec:p15-w02b-p15-029-variety}

\begin{hbdefinition}{def:p15-w02b-p15-029-variety}{The 110I requirement}
Two investigators share prior knowledge and each tests a law \(l\) on
one hundred specimens with all results positive; the first uses
specimens of one kind, the second of various kinds. Let \(e_1,e_2\) be
their total evidence. The requirement is that an adequate explicatum
\(c\) give \(l\), or the prediction \(h\) of a further instance, a
higher value on \(e_2\) than on \(e_1\)
\cite[sec.~110I, p.~575]{Carnap1950Probability}.
\end{hbdefinition}

Carnap asserts that \(c^*\) satisfies this requirement and refers for
the demonstration to his 1945 paper \emph{On Inductive Logic}, sec.~15;
the 1950 book contains no proof. The requirement is stated informally:
it does not say which sentences express ``kinds'' or which further
instance is meant. The following exact calculation shows how the answer
depends on that choice for \(c^*\); it is an independent illustration,
not a reconstruction of the 1945 argument.

\begin{hbtheorem}{thm:p15-w02b-p15-029-variety}{proved=annex:p15-w02b-p15-029-variety}{Variety under \(c^*\) in a two-predicate language}
Let the primitive predicates be \(K\) (a kind) and \(W\), so
\(\kappa=4\), and let \(l=\forall x\,W(x)\). Let \(e_1\) report one
hundred individuals all \(K\land W\), and \(e_2\) report fifty
\(K\land W\) and fifty \(\neg K\land W\). For a new individual \(a\),
\begin{enumerate}
 \item \(c^*_N(l,e_1)=c^*_N(l,e_2)\) for every \(N\geq100\), and
       \(c^*(W(a),e_1)=c^*(W(a),e_2)=51/52\);
 \item \(c^*(W(a),e_1\land\neg K(a))=1/2\), while
       \(c^*(W(a),e_2\land\neg K(a))=51/52\).
\end{enumerate}
\end{hbtheorem}

So under \(c^*\) the variety of the sample leaves the law and the
unqualified next instance unchanged, and raises the prediction for an
instance of the kind that the uniform sample omitted. Whether this is
what the requirement demands depends on its intended formalization.

\section{Source-to-host bridge for the 1950 items}
\label{sec:p15-w02b-p15-029-bridge}

The left column gives Carnap's notation or numbered item; the middle
column its form in this book; the right column whether Carnap printed a
proof (``printed'') or the proof cited here is an independent
reconstruction (``reconstructed''). Where Carnap prints a proof, the
reconstruction follows its route unless stated.

\begin{center}
\small
\begin{tabular}{@{}p{0.27\linewidth}p{0.42\linewidth}p{0.25\linewidth}@{}}
\hline
1950 item & Host form & Proof standing\\
\hline
\(\mathfrak L_N\), \(\mathfrak L_\infty\) (sec.~15A) &
\cref{def:p15-w02b-p15-029-systems}; \(\Omega_N\) of
\cref{def:p15-carnap-state-range} & definitions\\
in, pr, \(\pi\) & individual constants, primitive predicates, their
number (\(q\) in \cref{def:p15-carnap-state-range}) & definitions\\
state description \(\mathfrak Z\) (D18-1a) & element of \(\Omega_N\) &
definition\\
Str (D27-1), isomorphism (sec.~26) & renaming orbit,
\cref{def:p15-carnap-structure-width} & T27-2f used in the 110A proof\\
\(Q\)-properties, \(\kappa=2^\pi\) (T31-1) & \(Q_j\), \(k=2^q\) &
reconstructed: \cref{thm:p15-carnap-language-counts}\\
width \(w\), relative width \(w/\kappa\) (D32-1) & \(w(M)\), \(w/k\)
& definition\\
\(Q\)-numbers \(N_1,\ldots,N_\kappa\) (sec.~34) & counts
\((n_1,\ldots,n_k)\) & definition\\
\(\tau\) (number of Str in \(\mathfrak L_N\), T35-1d), \(\zeta_i\)
(states isomorphic to \(\mathfrak Z_i\)) & \(r_N\), \(\lvert O_z\rvert\) & reconstructed:
\cref{thm:p15-carnap-language-counts}\\
D55-1--4, regular \(m\), \(c\) & \cref{def:p15-carnap-regular-confirmation}
& definitions\\
D56-1, D56-2, T56-2, T56-4a & \cref{def:p15-w02b-p15-029-limits} &
T56-4a printed; others definitions\\
D81-1, T81-1 & \cref{def:p15-carnap-comparative,thm:p15-carnap-comparative}
& T81-1 printed; reconstruction follows the nine-sentence route\\
T81-2a & \cref{thm:p15-w02b-p15-029-t81-transfer} & printed, citing
T20-11; T20-11 step reconstructed in
\cref{thm:p15-w02b-p15-029-ltruth-transfer}\\
T94-1a, b & \cref{thm:p15-carnap-direct} & printed from T92-3;
reconstructed by counting\\
T95-1, T96-1 & \cref{thm:p15-w02b-p15-029-binomial,thm:p15-w02b-p15-029-bernoulli}
& T95-1 derivations indicated; T96-1 not printed; reconstructed\\
41B (1), (2); 41D (3)--(9) & \cref{def:p15-w02b-p15-029-betting,thm:p15-w02b-p15-029-betting}
& (8) from (7) printed; rest reconstructed\\
110A (1)(a)(b), (2) & \cref{def:p15-carnap-star},
\cref{thm:p15-carnap-star-construction} & existence printed (p.~564);
uniqueness reconstructed\\
110A (3) & \(m^\dagger\), \(c^\dagger\) of \cref{def:p15-carnap-star}
& example computation printed\\
110A (4) & mass formula of \cref{thm:p15-carnap-star-construction} &
substitution printed; reconstructed\\
110A (5) & \cref{thm:p15-w02b-p15-029-monadic-limit}(2) & not printed;
reconstructed\\
110B & direct inference, T94--96 & as above\\
110C (6), (7) & \cref{thm:p15-carnap-predictive} & not printed;
reconstructed\\
110D (8) & \cref{thm:p15-carnap-analogy} & not printed; reconstructed\\
110E (9) & inverse case of \cref{thm:p15-carnap-predictive} & not
printed; reconstructed\\
110F (10), (12), (13) & \cref{thm:p15-carnap-universal},
\cref{thm:p15-w02b-p15-029-monadic-limit}(3), and
\cref{thm:p15-w02b-p15-029-restricted-law} & not printed;
reconstructed\\
110F (11) & approximation for large \(s\) & not reconstructed\\
110G (14), (15), (17) & instance and qualified-instance values,
\cref{thm:p15-carnap-universal} & not printed; reconstructed\\
110I & \cref{def:p15-w02b-p15-029-variety} & proof referred to
Carnap 1945, sec.~15\\
\hline
\end{tabular}
\end{center}

\begin{hbcorollary}{thm:p15-w02b-p15-029-restricted-law}{proved=annex:p15-w02b-p15-029-restricted-law}{Formula (13): the restricted law}
Let \(M_1\) have width \(w_1\), \(0<w_1<\kappa\), let \(e'\) report
\(s\) individuals of which \(s_1\) are \(M_1\), and let \(l'\) say
that every individual of \(\mathfrak L_N\) outside the sample is
\(\neg M_1\). Then, for \(N\geq s\),
\[
 c^*_N(l',e')=
 \binom{s-s_1+N-s+\kappa-w_1-1}{N-s}\Big/\binom{N+\kappa-1}{N-s},
\]
which decreases strictly in \(s_1\) for fixed \(s,N\) with
\(N>s\), and for fixed \(s,s_1\) tends to zero as \(N\to\infty\).
\end{hbcorollary}

For \(s_1=0\) the restricted law is the law \(l\) itself and the
formula is (10), already proved as \cref{thm:p15-carnap-universal}.
The 1950 sources of the remaining numbered items, and the parts of
the book reconstructed in the two preceding chapters, are listed in
those chapters' boundaries.

\section{Zabell's history: probability, confirmation, exchangeability}
\label{sec:p15-w02b-p15-029-zabell}

Zabell's chapter writes Carnap's theory in event notation from a
Bayesian, subjective standpoint. Sentences are replaced by events, a
state description by an outcome sequence \(e_1,\ldots,e_n\) in
\(t\) categories, and a structure description by its counts
\(n_1,\ldots,n_t\) \cite[secs.~1, 4.2, pp.~265, 272]{Zabell2011CarnapInductiveInference}.
The held copy has the allocated pagination 265--309 and a 2009
copyright footer. The claims below are marked as historical,
philosophical, or mathematical.

\emph{Sec.~2.1 (historical).} Carnap introduced probability\(_1\) and
probability\(_2\) in 1945; the 1950 book assigns probabilities to
sentences; his later work uses events or propositions. Carnap's \(m\)
descends from the probability remarks of Wittgenstein and Waismann, and
\(c(h,e)=m(h\land e)/m(e)\). Hosiasson-Lindenbaum had studied
confirmation axiomatically: \(0\leq c\leq1\); invariance under
equivalence of either argument; \(c(h,e)=1\) if \(e\) implies \(h\);
additivity over hypotheses incompatible given \(e\); and
\(c(h\land h',e)=c(h,e)c(h',h\land e)\) (pp.~266--267).

\emph{Sec.~2.2 (historical and philosophical).} Ramsey derived utility
and probability from preferences; de Finetti took fair odds, with
coherence (no Dutch book) as the rationality constraint, and showed
that coherent odds obey the finitely additive axioms. Zabell credits
Kemeny (1955) with the converse: the five axioms are necessary and
sufficient for coherence. He reports that this persuaded Carnap to
adopt the betting interpretation, abandoning evidential support, one
of the three interpretations of 1950 (pp.~267--268). The 1950 betting
and estimate explanations are those of
\cref{thm:p15-w02b-p15-029-betting}; their mathematical content is not
Kemeny's theorem, whose proof is not reconstructed here.

\emph{Secs.~3--3.1 (mathematical and historical).} If trials are
conditionally independent with success probability \(p\) drawn from a
law \(\mu\) on \([0,1]\), the count \(S_n\) has the binomial mixture
law. The uniform \(\mu\) gives \(P(S_n=k)=1/(n+1)\); Bayes's
argument runs in the opposite direction, from that count law to the
uniform prior, by the uniqueness of a law on \([0,1]\) with given
moments; and the rule of succession follows (pp.~269--270). The
nineteenth-century criticism by Ellis, von Kries, Edgeworth, Venn,
Bertrand and Poincaré is reported as history.

\begin{hbtheorem}{thm:p15-w02b-p15-029-bayes-laplace}{proved=annex:p15-w02b-p15-029-bayes-laplace}{Bayes's postulate and the rule of succession}
For an infinite exchangeable binary sequence the following are
equivalent: (i) \(P(S_n=k)=1/(n+1)\) for all \(n\geq1\) and
\(0\leq k\leq n\); (ii) the de Finetti mixing law of
\cref{thm:p15-binary-definetti} is uniform on \([0,1]\). Under them,
\(P(X_{n+1}=1\mid S_n=k)=(k+1)/(n+2)\).
\end{hbtheorem}

\emph{Sec.~4 (mathematical and historical).} Johnson (1924) allows
\(t\geq2\) categories (4.1), introduces the permutation postulate,
which is exchangeability (4.2), and the combination postulate, giving
all count vectors of length \(n\) equal probability (4.3). Exchangeability
replaces the analogy with drawing from an urn of nature: instances are
like draws to the extent that they are judged exchangeable. Under it,
\(P(n_1,\ldots,n_t)=\frac{n!}{n_1!\cdots n_t!}P(e_1,\ldots,e_n)\)
and predictive probabilities depend only on the counts. With the
combination postulate one obtains \(m^*\) and \(c^*\) in the form
\((n_i+1)/(n+t)\) (pp.~271--274). The held copy prints the number of
count vectors as \(\binom{n+t-1}t\); the correct number is
\(\binom{n+t-1}{t-1}\) (for \(n=1\), \(t=2\) there are two), and
the formulas below use the correct value, which agrees with
\cref{thm:p15-carnap-language-counts}.

\begin{hbtheorem}{thm:p15-w02b-p15-029-johnson}{proved=annex:p15-w02b-p15-029-johnson}{Permutation and combination postulates}
Let \(X_1,\ldots,X_n\) take values in \(t\) categories with
exchangeable law. Then each count vector has probability
\(\frac{n!}{n_1!\cdots n_t!}\) times that of any sequence with those
counts, and \(P(X_{n+1}=c_i\mid X_1=e_1,\ldots,X_n=e_n)\) depends
only on the counts whenever \(X_1,\ldots,X_{n+1}\) is exchangeable and
the conditioning sequence has positive probability. If moreover all count
vectors have probability \(1/\binom{n+t-1}{t-1}\) at every length, a
sequence with counts \((n_i)\) has probability
\(\prod_in_i!\,(t-1)!/(n+t-1)!\), and the predictive probability is
\((n_i+1)/(n+t)\).
\end{hbtheorem}

\emph{Sec.~5 (mathematical and historical).} Johnson's sufficientness
postulate, so named by Good, requires the predictive probability of
category \(i\) to be a function \(f_i(n_i,n)\) of its own count and
the sample size only.

\begin{hbtheorem}{thm:p15-w02b-p15-029-johnson-carnap}{proved=annex:p15-w02b-p15-029-johnson-carnap}{The Johnson--Carnap continuum (Zabell, sec.~5.1)}
Let \(X_1,X_2,\ldots\) be an infinite exchangeable sequence with values
in \(t\geq2\) categories such that
\begin{enumerate}
 \item every finite sequence has positive probability;
 \item for every \(n\geq1\) there are functions \(f_i(\cdot,n)\) with
       \(P(X_{n+1}=c_i\mid X_1,\ldots,X_n)=f_i(n_i,n)\);
 \item either \(t\geq3\), or \(t=2\) and each \(f_i(\cdot,n)\) is
       affine: \(f_i(a,n)=g_i(n)+ad_i(n)\).
\end{enumerate}
Then either the \(X_j\) are independent with common law
\(P(X_1=c_i)=p_i>0\), or there are constants \(\alpha_i>0\), with
\(\alpha=\sum_i\alpha_i\), such that for all \(n\geq0\)
\[
 P(X_{n+1}=c_i\mid n_1,\ldots,n_t)=\frac{n_i+\alpha_i}{n+\alpha}.
\]
Conversely each such law satisfies (1)--(3)
(\cref{thm:p15-carnap-history-family}).
\end{hbtheorem}

Zabell states the theorem with conditions (1)--(3) for \(t\geq3\) and
says the case \(t=2\) is recovered when sufficientness is replaced by
linearity in \(n_i\), citing his 1982 paper; he notes that with two
categories sufficientness alone is vacuous (pp.~274--275). The proof
in the annex is an independent reconstruction along Johnson's
algebraic route; it uses infinite extendability only to exclude
negative \(\alpha_i\), which is why the finite continuum of his
sec.~8.4 admits negative parameters. Carnap's \(c^*\) is
\(\alpha_i=1\).

\emph{Sec.~5.2 (mathematical).} For an infinitely exchangeable
\(t\)-valued sequence there is a probability \(\mu\) on the simplex
\(\Delta_t\) with
\(P(n_1,\ldots,n_t)=\frac{n!}{n_1!\cdots n_t!}\int\prod_ip_i^{n_i}\,d\mu\)
for every \(n\); one \(\mu\) serves all \(n\), and no restriction is
placed on it (p.~275). This is the case of a finite state space in
\cref{thm:p15-definetti-kernel}, with \(\mu\) the mixing law; its
uniqueness is \cref{thm:p15-directing-law-uniqueness}. Zabell's
interpretation of the Johnson--Carnap theorem as identifying \(\mu\)
with a Dirichlet law, and Johnson's remark that his algebra replaces the
gamma-function integrals, are reported here; the Dirichlet integral
identity is not reconstructed in this chapter.

\emph{Secs.~5.3--5.4 (interpretive and historical).} For \(n>0\) the
continuum value is
\(\frac n{n+\alpha}\frac{n_i}n+\frac\alpha{n+\alpha}\frac{\alpha_i}\alpha\):
a weighted average of the observed frequency and the initial
probability \(\alpha_i/\alpha\), the weight on experience tending to one
as \(n\to\infty\) for every fixed \(\alpha\) (pp.~276--277). Johnson
treated the category-symmetric case \(\alpha_i=\alpha/t\); his proof,
published posthumously in 1932, had lacunae filled in Zabell's 1982
paper, which also treats the asymmetric case, \(t=\infty\), and finite
sequences. Carnap at first assumed both sufficientness and the
predictive form; Kemeny proved their equivalence for \(t>2\); Carnap
and Stegmüller treated \(t=2\); the category-asymmetric form appears
in Carnap's last work published by Jeffrey (1980) and in Kuipers (1978)
(p.~277).

\begin{proofboundary}{def:p15-w02b-p15-029-boundary}{What remains outside this chapter}
For the 1950 book: the approximation formula 110F (11), whose printed
form is not legible in the held text extraction; the proof of the 110I
assertion, which the book refers to Carnap 1945 (not held); and the
comparative vindication of \(c^*\) and a confirmation function for the
full quantitative language, which the 1950 text defers to a second
volume. Relational \(\mathfrak L_\infty\): limit existence for \(m^*\)
on general sentences is not asserted by Carnap outside the monadic
restriction of sec.~110 and is not established here.

For Zabell's chapter: Kemeny's coherence converse (sec.~2.2), the
Dirichlet identification (sec.~5.2), and the sections from 6 onward
other than the results already given in
Chapter~\ref{sec:p15-carnap-logical-probability-and-induction}: the
extended continuum confirming universal generalizations (6.2), the
Hintikka--Niiniluoto--Kuipers systems (6.3), the finite rule of
succession, finite continuum and quantitative finite approximation
(8.3--8.5), analogy, Markov and parametric exchangeability and group
invariance (10.1--10.5), species sampling and its three-parameter
characterization (11), the philosophical discussions of sec.~12
including ravens, grue, total evidence and Popper--Miller, and
secs.~13--14. The quantitative finite approximation is related to, but
not identical with, the bound in \cref{thm:p15-finite-urn-approximation}.
The held copy's 2009 footer against the 2011 publication record is not
resolved.
\end{proofboundary}
